A comunicação acadêmica nas universidades brasileiras acontece de forma
fragmentada, espalhada entre sistemas que não foram pensados para a interação
entre alunos, monitores e professores, o que acaba levando os estudantes a
recorrerem a canais informais, como o WhatsApp e o e-mail, nos quais o
conhecimento se perde com o tempo. Junto disso, as iniciativas de voluntariado
universitário dependem de redes sociais para a divulgação e de processos manuais
para registrar a participação e emitir os certificados. Este trabalho propõe o
desenvolvimento de uma plataforma web que junta, em um só lugar, um fórum
acadêmico organizado por disciplina e um sistema de voluntariado com inscrição e
certificação. Foi feito o levantamento bibliográfico sobre comunicação em
ambientes educacionais, fóruns de discussão, aprendizagem colaborativa e
voluntariado universitário, junto com a análise comparativa entre as plataformas
da UNIFEI, USP, UFMG e UNESPAR. Foram documentados 53 requisitos funcionais e
mais de 30 não funcionais, feita a modelagem do banco de dados com 17 tabelas em
terceira forma normal e definida a arquitetura do sistema, com o React.js, o
TypeScript e o TailwindCSS na camada de apresentação, o Python com o Django, o
Django REST Framework e o Django Channels na camada de aplicação, o PostgreSQL
como banco de dados relacional e o Redis para o gerenciamento de tokens e a
comunicação em tempo real. O escopo do produto mínimo apaixonante (MLP) foi
delimitado dando prioridade a funcionalidades como o login por CPF, a criação de
tópicos com votação e melhor resposta, a busca textual nos tópicos e a inscrição
em oportunidades de voluntariado com a certificação integrada. Para o TCC-2
estão previstas a implementação do protótipo e a validação com usuários da
UNIFEI.